% ==================== 第3章 恒定磁场 ====================
% 子文件：仅正文，由 1.电磁场笔记.tex \input 引入
% 教材：《工程电磁场导论》第3章 §3-1 ~ §3-9
% 公式统一用 formulaBox（本目录 3.矢量分析.tex 中定义）
% ==========================================================

\chapter{恒定磁场}

\textbf{研究对象}：恒定电流（不随时间变化）所激发的磁场，场量满足 $\partial(\bullet)/\partial t=0$。

\textbf{结构对照}：静电场由 $\nabla\times\vec{E}=0$ 引入标量电位、由 $\nabla\cdot\vec{D}=\rho$ 得泊松方程；恒定磁场恰好相反——由 $\nabla\cdot\vec{B}=0$ 引入\emph{矢量}磁位 $\vec{A}$、由 $\nabla\times\vec{H}=\vec{J}$ 得泊松方程。这个"叉与点互换"是本章公式的总来源。

% ---------------------------------------------------------
\section{磁感应强度}

\subsection{安培力定律与毕奥—沙伐定律}
真空中两细导线回路 $l$、$l'$（电流 $I$、$I'$）间的作用力为
\begin{formulaBox}
\[
\vec{F}=\frac{\mu_0}{4\pi}\oint_l\oint_{l'}\frac{I\,\mathrm{d}\vec{l}\times\big(I'\,\mathrm{d}\vec{l}'\times\vec{e}_R\big)}{R^2}
\]
\end{formulaBox}
$\mu_0=4\pi\times10^{-7}\ \mathrm{H/m}$ 为真空磁导率，$\vec{e}_R$ 由 $I'\mathrm{d}\vec{l}'$ 指向 $I\mathrm{d}\vec{l}$。它与库仑定律地位相当，但含双重矢量积，方向关系复杂——根源是电流元的矢量性；且只适用于\emph{闭合}回路（恒定电流必须闭合）。

把上式括号内的量理解为回路 $l'$ 在 $I\mathrm{d}\vec{l}$ 处产生的效应，定义\textbf{磁感应强度}（磁通密度）：
\begin{formulaBox}
\[
\vec{B}=\frac{\mu_0}{4\pi}\oint_{l'}\frac{I'\,\mathrm{d}\vec{l}'\times\vec{e}_R}{R^2}
=\frac{\mu_0}{4\pi}\int_{V'}\frac{\vec{J}\times\vec{e}_R}{R^2}\,\mathrm{d}V'
=\frac{\mu_0}{4\pi}\int_{S'}\frac{\vec{K}\times\vec{e}_R}{R^2}\,\mathrm{d}S'
\qquad(\mathrm{T})
\]
\end{formulaBox}
（依次对应线、体、面电流，即\textbf{毕奥—沙伐定律}。）于是安培力定律写成一般形式
\begin{formulaBox}
\[
\mathrm{d}\vec{f}=I\,\mathrm{d}\vec{l}\times\vec{B},
\qquad
\vec{F}=\oint_l I\,\mathrm{d}\vec{l}\times\vec{B}
\qquad(\text{均匀场中 }\vec{F}=I\vec{l}\times\vec{B})
\]
\end{formulaBox}
三个可直接引用的积分结论（$\rho$ 为场点到导线的垂距，$K$ 为面电流线密度）：
\begin{formulaBox}
\[
B_{\text{无限长直导线}}=\frac{\mu_0 I}{2\pi\rho},
\qquad
B_{\text{长 }2L\text{ 直导线中点垂距}\rho}=\frac{\mu_0 I L}{2\pi\rho\sqrt{\rho^2+L^2}},
\qquad
B_{\text{无限大电流片}}=\frac{\mu_0 K}{2}
\]
\end{formulaBox}

\subsection{洛伦兹力与磁感应线}
\begin{formulaBox}
\[
\vec{F}=q\vec{v}\times\vec{B},
\qquad
\vec{B}\times\mathrm{d}\vec{l}=0
\ \Big(\text{即}\ \frac{\mathrm{d}x}{B_x}=\frac{\mathrm{d}y}{B_y}=\frac{\mathrm{d}z}{B_z}\Big)
\]
\end{formulaBox}
前者为运动电荷受的\textbf{洛伦兹力}（静止电荷不受磁场力；力恒垂直于速度，只改方向不改量值，故\textbf{洛伦兹力不做功}）；后者是\textbf{磁感应线}的微分方程。

\subsection{磁通与磁通连续性原理}
\begin{formulaBox}
\[
\Phi=\iint_S \vec{B}\cdot\mathrm{d}\vec{S}\ (\mathrm{Wb}),
\qquad
\oint_S \vec{B}\cdot\mathrm{d}\vec{S}=0
\qquad\Longleftrightarrow\qquad
\nabla\cdot\vec{B}=0
\]
\end{formulaBox}
磁感应线闭合、无始无终，自然界不存在与电荷对应的磁荷，故穿出任意闭合面的净磁通为零——\textbf{磁通连续性原理}：恒定磁场是\emph{无散场}，穿出等于穿入，同一闭合曲线为周界的各曲面上磁通相同。

\corepoint{$\nabla\cdot\vec{B}=0$ 只说明"无磁荷源"，不说明磁场没有源——磁场的涡旋源是电流。另注意 $\Phi$ 是穿过某个\emph{面}的通量，$\oint_S\vec{B}\cdot\mathrm{d}\vec{S}=0$ 只对\emph{闭合}面成立。}

% ---------------------------------------------------------
\section{安培环路定律}

\subsection{真空中的安培环路定律}
\begin{formulaBox}
\[
\oint_l \vec{B}\cdot\mathrm{d}\vec{l}=\mu_0\sum I_i
\qquad\Longleftrightarrow\qquad
\nabla\times\vec{B}=\mu_0\vec{J}_c
\]
\end{formulaBox}
电流方向与回路绕向符合\textbf{右手螺旋关系}时取正。磁场是\textbf{有旋场}，$\vec{B}$ 的环流一般不恒为零，故磁场非保守场、一般无"磁势"；磁感线与电流垂直并相互交链。

\textbf{用它求 $\vec{B}$ 的条件}：磁场具有某种对称性，且能选到这样的积分路径——路径上各点 $\vec{B}$ 与 $\mathrm{d}\vec{l}$ 夹角相同、$B$ 值相等，从而把 $\oint$ 化为代数式。典型应用如载流圆柱（$\rho<R$ 时 $B=\mu_0I\rho/(2\pi R^2)$，$\rho>R$ 时 $B=\mu_0I/(2\pi\rho)$）、同轴电缆、螺环等。

\subsection{媒质的磁化与磁化强度}
分子（原子）中电子的轨道运动与自旋都等效于电流，其对外磁效应的总可用一个等效圆电流表示，称\textbf{分子电流}（束缚电流），它不引起电荷迁移，但能激发磁场。分子电流具有\textbf{分子磁矩} $\vec{m}=I\vec{S}$（$\vec{S}$ 方向与电流绕向服从右手螺旋）。无外场时分子磁矩因热运动随机排列、总磁矩为零；放入外场后受转矩 $\vec{T}=\vec{m}\times\vec{B}$ 而趋于与外场同向，物质显磁性——即\textbf{磁化}，其宏观后果是出现\textbf{磁化电流}，附加磁场改变原场分布。

\textbf{磁化强度}（单位体积内分子磁矩的矢量和）与磁化电流的关系为
\begin{formulaBox}
\[
\vec{M}=\lim_{\Delta V\to 0}\frac{\sum_i\vec{m}_i}{\Delta V}\ (\mathrm{A/m}),
\qquad
\vec{J}_m=\nabla\times\vec{M},
\qquad
\oint_l\vec{M}\cdot\mathrm{d}\vec{l}=I_m,
\qquad
\vec{K}_m=\vec{M}\times\vec{e}_n
\]
\end{formulaBox}
于是"有磁介质时的磁场"可化为"把磁化电流与自由电流一起在真空中求场"。

\subsection{磁场强度与一般形式的安培环路定律}
由 $\nabla\times\vec{B}=\mu_0(\vec{J}_c+\vec{J}_m)$、$\vec{J}_m=\nabla\times\vec{M}$ 移项得
\begin{formulaBox}
\[
\nabla\times\Big(\frac{\vec{B}}{\mu_0}-\vec{M}\Big)=\vec{J}_c
\ \Longrightarrow\
\vec{H}=\frac{\vec{B}}{\mu_0}-\vec{M},
\qquad
\nabla\times\vec{H}=\vec{J}_c
\ \Longleftrightarrow\
\oint_l\vec{H}\cdot\mathrm{d}\vec{l}=\sum I
\]
\end{formulaBox}
$\vec{H}$ 称\textbf{磁场强度}（$\mathrm{A/m}$）。\textbf{物理意义}：$\vec{H}$ 的\emph{环流}只由自由电流决定、与磁化电流（即媒质分布）无关；但 $\vec{H}$ 的\emph{分布}仍与媒质有关，两者不可混同。

\subsection{磁介质的物性方程}
\begin{formulaBox}
\[
\vec{B}=\mu_0(\vec{H}+\vec{M}),
\quad
\vec{M}=\chi_m\vec{H},
\quad
\vec{B}=\mu_0(1+\chi_m)\vec{H}=\mu_0\mu_r\vec{H}=\mu\vec{H}
\]
\end{formulaBox}
$\chi_m$ 为磁化率，$\mu_r=1+\chi_m$ 为相对磁导率，$\mu=\mu_0\mu_r$ 为磁导率（$\mathrm{H/m}$）。顺磁质 $\chi_m>0$、抗磁质 $\chi_m<0$（均极小），铁磁质 $\mu_r$ 可达 $10^3\sim10^5$ 且非线性、有磁滞与饱和。注意 $\vec{B}=\mu\vec{H}$ \textbf{仅适用于各向同性线性}媒质，而 $\vec{B}=\mu_0(\vec{H}+\vec{M})$ 无此限制。

\textbf{推论}：若电流周围无限充满均匀导磁媒质，则各点 $\vec{B}$ 方向同真空、量值增大 $\mu_r$ 倍（把 $\mu_0$ 换成 $\mu$）。

\subsection{电介质极化与磁介质磁化的对照}
\begin{center}
\small
\begin{tabularx}{\linewidth}{@{}l X X@{}}
\toprule
 & \textbf{电介质极化} & \textbf{磁介质磁化} \\
\midrule
基本单元 & 电偶极矩 $\vec{p}=q\vec{d}$ & 分子磁矩 $\vec{m}=I\vec{S}$ \\
强度量 & $\vec{P}=\lim\dfrac{\sum\vec{p}_i}{\Delta V}$ & $\vec{M}=\lim\dfrac{\sum\vec{m}_i}{\Delta V}$ \\
附加源 & $\sigma_p=\vec{P}\cdot\vec{e}_n$，$\rho_p=-\nabla\cdot\vec{P}$ & $\vec{K}_m=\vec{M}\times\vec{e}_n$，$\vec{J}_m=\nabla\times\vec{M}$ \\
辅助场量 & $\vec{D}=\varepsilon_0\vec{E}+\vec{P}$ & $\vec{H}=\vec{B}/\mu_0-\vec{M}$ \\
媒质方程 & $\vec{D}=\varepsilon\vec{E}$ & $\vec{B}=\mu\vec{H}$ \\
散/旋 & $\nabla\cdot\vec{D}=\rho$，$\nabla\times\vec{E}=0$ & $\nabla\cdot\vec{B}=0$，$\nabla\times\vec{H}=\vec{J}_c$ \\
\bottomrule
\end{tabularx}
\end{center}

\corepoint{最易写错的两组关系：\emph{电}——$\vec{D}$ 的散度由自由电荷决定、$\vec{E}$ 无旋；\emph{磁}——$\vec{B}$ 的散度恒为零、$\vec{H}$ 的旋度由自由电流决定。}

% ---------------------------------------------------------
\section{恒定磁场的基本方程与分界面上的衔接条件}

\subsection{恒定磁场的基本方程组}
\begin{formulaBox}
\[
\begin{cases}
\displaystyle\oint_S\vec{B}\cdot\mathrm{d}\vec{S}=0\\[6pt]
\displaystyle\oint_l\vec{H}\cdot\mathrm{d}\vec{l}=\sum I
\end{cases}
\qquad
\begin{cases}
\nabla\cdot\vec{B}=0\\[4pt]
\nabla\times\vec{H}=\vec{J}_c
\end{cases}
\qquad
\vec{B}=\mu\vec{H}
\]
\end{formulaBox}
$\nabla\cdot\vec{B}=0$：无源场，磁感应线无头无尾、无磁荷；$\nabla\times\vec{H}=\vec{J}_c$：有旋场，电流密度是涡旋源；$\vec{B}=\mu\vec{H}$：各向同性线性媒质的构成方程（铁磁质不成立）。

\subsection{分界面上的衔接条件}
设 $\vec{e}_n$ 由媒质 1 指向媒质 2：
\begin{formulaBox}
\[
\begin{cases}
H_{1t}-H_{2t}=K\ \ \text{或}\ \ \vec{e}_n\times(\vec{H}_2-\vec{H}_1)=\vec{K} & \text{（切向：}H_t\text{跃变）}\\[4pt]
B_{1n}=B_{2n}\ \ \text{或}\ \ \vec{e}_n\cdot(\vec{B}_2-\vec{B}_1)=0 & \text{（法向：}B_n\text{连续）}
\end{cases}
\]
\end{formulaBox}
$H_t$ 的跃变量等于界面上的\textbf{面电流线密度} $K$（由跨界面小矩形回路的 $\oint\vec{H}\cdot\mathrm{d}\vec{l}=\sum I$ 得到）；$B_n$ 连续由跨界面小扁盒的 $\oint\vec{B}\cdot\mathrm{d}\vec{S}=0$ 得到。当界面上无面电流（$K=0$）时 $H_t$ 也连续，但 $B_t$ 仍不连续（$B_t=\mu H_t$，两侧 $\mu$ 不同）。理论上只有理想导体、超导体表面才可能有面电流；工程上仅当集肤效应十分显著时才可近似视为面电流。

\subsection{磁场折射定律}
分界面上无面电流、两侧均为各向同性线性媒质时
\begin{formulaBox}
\[
\frac{\tan\theta_1}{\tan\theta_2}=\frac{\mu_1}{\mu_2}
\]
\end{formulaBox}
若媒质 2 为铁磁质（$\mu_2\gg\mu_1$），除 $\theta_1=90^\circ$ 外 $\theta_2\to90^\circ$：铁磁质内 $B$ 线几乎平行于界面，且 $B_2\gg B_1$，磁通高度集中。即"\textbf{$B$ 线沿铁走}"，这是磁路概念的物理基础。

% ---------------------------------------------------------
\section{磁矢位·恒定磁场的边值问题}

\subsection{磁矢位的引出与库仑规范}
因 $\nabla\cdot\vec{B}=0$，可引入矢量函数使
\begin{formulaBox}
\[
\vec{B}=\nabla\times\vec{A}\qquad(\mathrm{Wb/m})
\]
\end{formulaBox}
代入 $\nabla\times\vec{H}=\vec{J}$、$\vec{B}=\mu\vec{H}$，用 $\nabla\times\nabla\times\vec{A}=\nabla(\nabla\cdot\vec{A})-\nabla^2\vec{A}$：
\begin{formulaBox}
\[
\nabla(\nabla\cdot\vec{A})-\nabla^2\vec{A}=-\mu\vec{J}
\]
\end{formulaBox}
唯一确定矢量场须同时给定旋度与散度，故再规定
\begin{formulaBox}
\[
\nabla\cdot\vec{A}=0\ \text{（库仑规范）}
\qquad\Longrightarrow\qquad
\nabla^2\vec{A}=-\mu\vec{J}\ \text{（矢量泊松方程）}
\]
\end{formulaBox}
直角坐标下即三个形式相同的标量泊松方程 $\nabla^2A_x=-\mu J_x$ 等。当电流分布于有限空间、取无限远处 $\vec{A}=0$ 时
\begin{formulaBox}
\[
\vec{A}=\frac{\mu_0}{4\pi}\int_{V'}\frac{\vec{J}}{R}\mathrm{d}V'
=\frac{\mu_0}{4\pi}\int_{S'}\frac{\vec{K}}{R}\mathrm{d}S'
=\frac{\mu_0}{4\pi}\oint_{l'}\frac{I\,\mathrm{d}\vec{l}'}{R}
\]
\end{formulaBox}
每个元电流产生的 $\vec{A}$ 与该元电流\emph{同方向}（与标量电位的根本差别），故可分量计算；$\vec{A}$ 无物理意义，但便于求 $\vec{B}$、磁通与电感。

\subsection{磁矢位与磁通}
由斯托克斯定理
\begin{formulaBox}
\[
\Phi=\iint_S\vec{B}\cdot\mathrm{d}\vec{S}=\iint_S(\nabla\times\vec{A})\cdot\mathrm{d}\vec{S}=\oint_l\vec{A}\cdot\mathrm{d}\vec{l}
\]
\end{formulaBox}
即 $\vec{A}$ 沿闭合回路的环量等于穿过以该回路为周界的曲面的磁通。对\textbf{平行平面磁场}（$\vec{A}=A_z\vec{e}_z$），$B$ 线即等 $A$ 线，等 $A$ 线密处 $B$ 大。

\subsection{磁矢位的衔接条件与边值问题}
由跨界面小矩形回路的 $\Phi=\oint\vec{A}\cdot\mathrm{d}\vec{l}$（高度趋于零）得切向分量连续，由 $\nabla\cdot\vec{A}=0$（跨界面小扁圆柱）得法向分量连续，故
\begin{formulaBox}
\[
\vec{A}_1=\vec{A}_2\ \text{（界面上 $\vec{A}$ 连续）}
\]
\end{formulaBox}
另一条件由 $\vec{e}_n\times(\vec{H}_2-\vec{H}_1)=\vec{K}$ 换算。对平行平面磁场简化为
\begin{formulaBox}
\[
A_1=A_2,
\qquad
\frac{1}{\mu_1}\frac{\partial A_1}{\partial n}-\frac{1}{\mu_2}\frac{\partial A_2}{\partial n}=K
\]
\end{formulaBox}
"$\nabla^2\vec{A}=-\mu\vec{J}$（无电流处 $\nabla^2\vec{A}=0$）$+$ 衔接条件 $+$ 场域边界条件"即恒定磁场的\textbf{边值问题}（磁矢位表述）。

\subsection{标量磁位}
在\textbf{传导电流处处为零}的区域（$\vec{J}_c=0$），$\nabla\times\vec{H}=0$，可有条件引入\textbf{标量磁位}：
\begin{formulaBox}
\[
\vec{H}=-\nabla\varphi_m,
\quad
\varphi_{mP}=\int_P^Q\vec{H}\cdot\mathrm{d}\vec{l}\ (\mathrm{A}),
\quad
U_{mAB}=\int_A^B\vec{H}\cdot\mathrm{d}\vec{l}=\varphi_{mA}-\varphi_{mB}
\]
\end{formulaBox}
均匀媒质中 $\nabla^2\varphi_m=0$（拉普拉斯方程）；衔接条件为
\begin{formulaBox}
\[
\varphi_{m1}=\varphi_{m2},
\qquad
\mu_1\frac{\partial\varphi_{m1}}{\partial n}=\mu_2\frac{\partial\varphi_{m2}}{\partial n}
\]
\end{formulaBox}
\textbf{多值性}：两点间磁压与积分路径有关，若回路绕电流 $k$ 次则差 $kI$，故即使选定参考点 $\varphi_m$ 仍是多值函数。消除办法是规定积分路径不穿过电流回路所限定的面（\textbf{磁屏障面}），使磁位单值。多值性不影响 $\vec{B}$、$\vec{H}$ 的计算。

\textbf{与电位的三点区别}：(1) 存在条件不同——$\varphi_m$ 只在 $\vec{J}_c=0$ 区可引入；(2) 多值性来源不同——电位只由参考点决定，磁位还受路径影响；(3) 物理意义不同——$\varphi_m$ 纯为计算方便，无物理意义。另：因 $H_t\approx0$，磁导率很大的材料表面近似为\textbf{等磁位面}，此结论在磁路计算中反复使用。

\corepoint{$\vec{A}$ 与 $\varphi_m$ 的分工：有电流的区域只能用 $\vec{A}$（矢量泊松方程），无电流区域才用 $\varphi_m$（标量拉普拉斯方程）。解题第一步先判断区域内有无电流，选错位函数则方程不成立。}

% ---------------------------------------------------------
\section{镜像法}

\subsection{依据}
由\textbf{唯一性定理}，若所求区域内方程相同、边界条件相同，解唯一。故可用镜像电流代替分界面上的（磁化）束缚电流，把分区问题化为均匀空间的场叠加——即\textbf{磁镜像法}。

\subsection{线电流对无限大媒质分界面的镜像}
媒质 1、2 的磁导率为 $\mu_1$、$\mu_2$，媒质 1 中有平行于界面的无限长直导线（电流 $I$）。求媒质 1 的场时，认为全空间充满 $\mu_1$，场由 $I$ 与对称位置的像电流 $I'$ 共同产生；求媒质 2 的场时，认为全空间充满 $\mu_2$，场由原位置的像电流 $I''$ 单独产生。由 $H_{2t}=H_{1t}$、$B_{1n}=B_{2n}$ 联立得
\begin{formulaBox}
\[
I'=\frac{\mu_2-\mu_1}{\mu_2+\mu_1}I,
\qquad
I''=\frac{2\mu_1}{\mu_2+\mu_1}I
\]
\end{formulaBox}
$I'$、$I''$ 的\textbf{参考方向均规定与 $I$ 相同}（正负由此体现）；$I''$ 恒为正，$I'$ 的正负取决于 $(\mu_2-\mu_1)$。用法记忆：\emph{媒质 1 场} $=I+I'$（全在 $\mu_1$ 中）；\emph{媒质 2 场} $=I''$（全在 $\mu_2$ 中），两者不可混用于同一区域。

\subsection{两种典型特例}
\begin{enumerate}
  \item \textbf{导线在空气中，另一侧为铁磁质}（$\mu_1=\mu_0$，$\mu_2\to\infty$）：$I'\approx I$，$I''\approx0$。铁磁质内 $H_2\approx0$，但\emph{不能}认为 $B_2$ 也为零——由 $B_n$ 连续，磁通仍穿入铁磁质。
  \item \textbf{导线在铁磁质中，另一侧为空气}（$\mu_1\to\infty$，$\mu_2=\mu_0$）：$I'\approx-I$，$I''\approx2I$。空气中的磁感应强度比全空间为空气时\textbf{增大一倍}（电流相同）。
\end{enumerate}

\subsection{适用条件与易错点}
\begin{enumerate}
  \item 界面须为无限大平面，两媒质各自均匀、各向同性、线性；
  \item 镜像电流位置由对称确定，符号约定须统一（都相对 $I$ 的参考方向）；
  \item 只适用于线电流及可等效为线电流的情形，有限长导线与线圈需另行处理。
\end{enumerate}

\corepoint{磁镜像与电镜像"符号相反"：静电场中电荷靠近理想导体（$\varepsilon_2\to\infty$）时镜像电荷异号（$-q$）；磁场中电流靠近铁磁质（$\mu_2\to\infty$）时镜像电流同号（$+I$）。记不住就回到 $B_n$ 连续这一条件重推——铁磁质"导通磁通"，导体"终止电力线"，性质相反。}

% ---------------------------------------------------------
\section{电感}

\subsection{自感}
线性媒质中磁链与电流成正比：
\begin{formulaBox}
\[
\psi=LI,
\qquad
L=\frac{\psi}{I}\ (\mathrm{H}),
\qquad
L=L_i+L_o
\]
\end{formulaBox}
$L$ 只与回路尺寸、几何形状及媒质分布有关，与电流、磁链数值无关。
\begin{itemize}
  \item \textbf{内自感} $L_i=\psi_i/I$：内磁通在导线内部闭合，只与\emph{部分}电流交链，求元内磁链时交链电流按 $(I'/I)=(\rho/R)^2$ 折算，故 $\mathrm{d}\psi_i=(\rho^2/R^2)\mathrm{d}\Phi$；
  \item \textbf{外自感} $L_o=\psi_o/I$：外磁通在导线外部闭合，与全部电流交链，$\mathrm{d}\psi_o=\mathrm{d}\Phi$。
\end{itemize}
圆截面导线的内自感为
\begin{formulaBox}
\[
L_i=\frac{\mu_0 l}{8\pi}
\qquad\text{（只与长度 }l\text{ 有关，与半径无关）}
\]
\end{formulaBox}

\subsection{两个必须掌握的自感结果}
\textbf{同轴电缆}（长 $l$，内导体半径 $R_1$，外导体内半径 $R_2$，外壳很薄，磁导率均为 $\mu_0$）：
\begin{formulaBox}
\[
L=\frac{\mu_0 l}{8\pi}+\frac{\mu_0 l}{2\pi}\ln\frac{R_2}{R_1}
\]
\end{formulaBox}
\textbf{二线传输线}（半径 $R$ 的两长直导线，轴线间距 $D\gg R$）：
\begin{formulaBox}
\[
L=\frac{\mu_0 l}{\pi}\Big[\ln\frac{D}{R}+\frac14\Big]
\]
\end{formulaBox}
（$\ln(D/R)$ 为外自感，$1/4$ 为两根导线内自感之和。）

\textbf{计算步骤}：设电流 $I\to$ 由安培环路定律求 $\vec{H}$、$\vec{B}=\mu\vec{H}\to$ 取面元求 $\mathrm{d}\Phi=B\mathrm{d}S\to$ 按交链电流比例得 $\mathrm{d}\psi$、积分得 $\psi=\psi_i+\psi_o\to L=\psi/I$。

\subsection{互感}
\begin{formulaBox}
\[
\psi_{21}=M_{21}I_1,
\qquad
M_{21}=\frac{\psi_{21}}{I_1},
\qquad
M_{12}=M_{21}=M\ (\mathrm{H})
\]
\end{formulaBox}
互感不仅与两回路的形状、尺寸、磁导率有关，还与\textbf{相互位置}有关。对细导线回路的一般公式（\textbf{诺伊曼公式}）为
\begin{formulaBox}
\[
M=\frac{\mu_0}{4\pi}\oint_{l_1}\oint_{l_2}\frac{\mathrm{d}\vec{l}_1\cdot\mathrm{d}\vec{l}_2}{R}
\qquad
\Big(\text{紧绕 }N_1,N_2\text{ 匝时乘 }N_1N_2\Big)
\]
\end{formulaBox}
它说明互感只取决于两回路的几何，无需先求场；但双重线积分通常较繁，工程上多用磁链法直接算 $\psi_{21}$。

\subsection{用磁场能量求电感}
对单个载流回路 $W_m=\frac12 LI^2$，故
\begin{formulaBox}
\[
L=\frac{2W_m}{I^2}
\]
\end{formulaBox}
当 $\vec{B}$、$\vec{H}$ 已用数值方法求出时，能量法求电感比磁链法更方便，是工程数值计算的常用做法。

% ---------------------------------------------------------
\section{磁场能量与力}

\subsection{磁场能量}
设磁场建立过程缓慢、媒质线性、无辐射与其它损耗，则外源做功全部转为磁场能量：
\begin{formulaBox}
\[
W_m=\frac12\sum_{k=1}^{n}I_k\psi_k,
\qquad
W_m=\frac12\int_V\vec{A}\cdot\vec{J}\,\mathrm{d}V
=\frac12\int_V\vec{B}\cdot\vec{H}\,\mathrm{d}V
\]
\end{formulaBox}
单回路为 $W_m=\frac12LI^2$；两回路为
\begin{formulaBox}
\[
W_m=\frac12L_1I_1^2+\frac12L_2I_2^2+MI_1I_2
\]
\end{formulaBox}
前两项为\textbf{自有能}（恒正，只与各自回路的电流、自感有关）；$MI_1I_2$ 为\textbf{互有能}（可正可负：两回路电流同时从同名端流入时取正）。可见磁场能量只与最终状态有关，与建立过程无关。

\subsection{磁场能量密度}
\begin{formulaBox}
\[
w_m=\frac12\vec{B}\cdot\vec{H}=\frac{B^2}{2\mu}=\frac12\mu H^2
\qquad(\mathrm{J/m^3})
\]
\end{formulaBox}
磁场能量并非只存在于电流回路内，而是\textbf{分布在整个磁场空间}——场是能量的载体，与静电场 $w_e=\frac12\vec{D}\cdot\vec{E}$ 完全对应。

\subsection{虚位移法求磁场力}
设 $n$ 个载流回路接于电压源，仅第 $p$ 号回路可沿广义坐标 $g$ 移动。能量守恒为
\begin{formulaBox}
\[
\mathrm{d}W=\mathrm{d}W_m+f\,\mathrm{d}g,
\qquad
\mathrm{d}W=\sum_k I_k\,\mathrm{d}\psi_k
\]
\end{formulaBox}
\begin{enumerate}
  \item \textbf{电流不变}（各回路接恒流源）：$\mathrm{d}W=2\mathrm{d}W_m$，外源供能一半储能、一半作机械功，
  \begin{formulaBox}
  \[
  f=\left.\frac{\partial W_m}{\partial g}\right|_{I=\text{常数}}
  \]
  \end{formulaBox}
  \item \textbf{磁链不变}（回路开路、无感应电流）：$\mathrm{d}W=0$，机械功只能取自磁场能量的减少，
  \begin{formulaBox}
  \[
  f=-\left.\frac{\partial W_m}{\partial g}\right|_{\psi=\text{常数}}
  \]
  \end{formulaBox}
\end{enumerate}
两式在当时状态下给出同一个力，可任选其一。求相互作用力时，写出相互能 $MI_1I_2$ 后对 $g$ 求偏导即可。

\textbf{电磁铁起重力}：气隙磁场近似均匀、铁芯内储能可忽略，每个气隙储能 $W_m=\dfrac{B^2}{2\mu_0}Sl$，对气隙长度 $l$ 求偏导：
\begin{formulaBox}
\[
F=\frac{B^2S}{2\mu_0}
\qquad\text{（负号表示力使气隙缩短，表现为吸力）}
\]
\end{formulaBox}

\subsection{磁场力面密度与法拉第观点}
\begin{formulaBox}
\[
f_0=\frac12\frac{B^2}{\mu_0}=\frac12\mu_0H^2
\qquad\text{（等于该处磁场能量体密度）}
\]
\end{formulaBox}
\textbf{法拉第观点}：每束磁感应线形成的"磁感应管"沿轴向受\textbf{纵张力}、垂直方向受\textbf{侧压力}，二者数值相等：
\begin{formulaBox}
\[
\frac12BH=\frac12\mu H^2=\frac{B^2}{2\mu}
\]
\end{formulaBox}
据此可定性判断受力（气隙中 $B$ 管沿轴向有收缩趋势，故磁极间表现为吸力）。对两种媒质的分界面：
\begin{formulaBox}
\[
f_0=\frac{\mu_2-\mu_1}{2\mu_1\mu_2}\big(B_{1n}^2+\mu_1\mu_2H_{1t}^2\big)
\]
\end{formulaBox}
无论磁场方向如何，此力\textbf{总垂直于界面}，且\textbf{总是由磁导率较大的媒质指向较小的媒质}。

\corepoint{虚位移法三处高频错误：(1) 恒流与恒磁链的符号用反（恒流取 $+\partial W_m/\partial g$，恒磁链取 $-\partial W_m/\partial g$）；(2) 忘记恒流时"外源供能是储能增量的两倍"；(3) 把 $f_0=B^2/(2\mu_0)$ 当总力——它只是单位面积上的力。}

% ---------------------------------------------------------
\section{磁路及其计算}

\subsection{磁路的概念}
铁磁材料 $\mu\gg\mu_0$，由折射定律知 $B$ 线几乎平行于铁磁质表面、并高度集中于内部（"$B$ 线沿铁走"）。把铁磁质做成闭合或基本闭合的形状，磁通便主要集中在铁心内，与电流集中在导体内相似。这种\textbf{能使磁通集中通过的路径}称为\textbf{磁路}。
\begin{itemize}
  \item \textbf{主磁通} $\Phi$：沿磁路（含气隙）闭合的磁通；
  \item \textbf{漏磁通}：穿出铁心、经磁路周围非铁磁质闭合的磁通。
\end{itemize}
绕在铁心上的线圈通以较小的励磁电流即可获得较强磁场，且磁场差不多被约束在磁路内，周围非铁磁质中很弱。

\subsection{磁路定律}
对无分支闭合铁心，沿一条闭合磁力线应用安培环路定律，并由 $\Phi=BS$、$H=\Phi/(\mu S)$ 得
\begin{formulaBox}
\[
\oint_l\vec{H}\cdot\mathrm{d}\vec{l}=NI
\ \Longrightarrow\
\Phi\int_l\frac{\mathrm{d}l}{\mu S}=NI
\]
\end{formulaBox}
与电阻 $R=\int\mathrm{d}l/(\gamma S)$ 对比，定义\textbf{磁阻}与\textbf{磁动势}
\begin{formulaBox}
\[
R_m=\int_l\frac{\mathrm{d}l}{\mu S}
\ \ \Big(\text{均匀段 }R_m=\frac{l}{\mu S}\Big),
\qquad
F_m=NI\ (\mathrm{A})
\]
\end{formulaBox}
得\textbf{无分支闭合磁路的欧姆定律}
\begin{formulaBox}
\[
\Phi R_m=F_m
\]
\end{formulaBox}
对有分支磁路，由磁通连续性与安培环路定律得\textbf{磁路定律}（对应电路的两个基尔霍夫定律）：
\begin{formulaBox}
\[
\sum\Phi=0\ (\text{节点}),
\qquad
\sum\Phi R_m=\sum F_m\ (\text{回路})
\]
\end{formulaBox}
磁阻的串并联与电阻同形：\textbf{串联}（各段 $\Phi$ 相同）$R_m=\sum_k R_{mk}$，铁心带气隙即两段磁阻串联；\textbf{并联}（各分支磁压相同）$\dfrac{1}{R_m}=\sum_k\dfrac{1}{R_{mk}}$。
\begin{center}
\small
\begin{tabularx}{\linewidth}{@{}l X X@{}}
\toprule
 & \textbf{电路} & \textbf{磁路} \\
\midrule
"流量" & 电流 $I$（连续） & 磁通 $\Phi$（连续） \\
"势" & 电动势 $\varepsilon$（电源） & 磁动势 $F_m=NI$（励磁线圈） \\
"阻力" & 电阻 $R=\dfrac{l}{\gamma S}$ & 磁阻 $R_m=\dfrac{l}{\mu S}$ \\
物性参数 & 电导率 $\gamma$ & 磁导率 $\mu$ \\
定律 & $IR=\varepsilon$，$\sum I=0$，$\sum IR=\sum\varepsilon$ & $\Phi R_m=F_m$，$\sum\Phi=0$，$\sum\Phi R_m=\sum F_m$ \\
"绝缘" & 理想绝缘体 $\gamma=0$，无漏电流 & \textbf{不存在磁绝缘体}，总有漏磁 \\
\bottomrule
\end{tabularx}
\end{center}
磁路定律是在"不计漏磁、$B$ 线沿铁心走向、截面上 $B$ 均匀"等近似下得到的，本质是\textbf{估算}，但对电机、变压器、电磁铁与仪表设计十分必要（含永磁体的磁路更复杂，此处从略）。

由 $\psi=N\Phi$ 还可得电感与磁阻的关系
\begin{formulaBox}
\[
L=\frac{\psi}{I}=\frac{N\Phi}{I}=\frac{N^2}{R_m}
\]
\end{formulaBox}
即\textbf{磁阻越小、电感越大}——这解释了铁心线圈为何电感远大于空心线圈，也解释了开气隙后电感为何急剧下降。

\subsection{磁路计算}
恒定磁通磁路（$\Phi$ 不随时间变化）的计算目的是找出 $\Phi$ 与 $F_m=NI$ 的关系，分两类问题：
\begin{enumerate}
  \item \textbf{已知 $\Phi$ 求 $NI$}（正面问题）：$\Phi$ 与各段截面 $\to$ 各段 $B$；查材料的 $B$—$H$ 曲线得各段 $H$（线性段可直接 $H=B/\mu$）；求各段磁压 $Hl$；相加得 $NI=\sum Hl$。
  \item \textbf{已知 $NI$ 求 $\Phi$}（反面问题）：铁磁质非线性使 $\mu$ 随 $H$ 变化，不能一次算出，须用\textbf{试探法}——先假设 $\Phi$（例如粗略认为全部磁动势都降在气隙上），算出对应 $NI$ 与给定值比较后修正，反复数次直至吻合；实质是"已知 $\Phi$ 求 $NI$"的多次重复。
\end{enumerate}
\textbf{气隙的影响}：$R_\delta=\dfrac{\delta}{\mu_0 S}$，因 $\mu_0\ll\mu_{\text{铁}}$，即使 $\delta$ 很小，气隙磁压降也常占全部磁动势的绝大部分——\emph{气隙决定磁路工作状态}，这是磁路计算第一条经验。

\textbf{边缘效应}：气隙处 $B$ 线向外扩散，实际有效面积略大于铁心截面，气隙磁阻比 $\delta/(\mu_0S)$ 略小。工程上用"有效面积"（如按 $(a+\delta)(b+\delta)$ 放大气隙截面）或大于 1 的卡氏系数修正，$\delta$ 越小越可忽略。

\subsection{磁屏蔽}
用高 $\mu$ 铁磁质做屏蔽罩，把需保护的空间罩起来。用磁阻并联解释：外界磁通遇到"铁壳壁"与"空腔"两条并联路径，空腔的 $\mu$ 远小于罩壁、磁阻远大于罩壁，故绝大部分 $B$ 线沿罩壁通过，进入空腔的很少。

\begin{itemize}
  \item 与静电屏蔽不同：闭合导体空腔内静电场为零，而外磁场中闭合铁磁空腔内的磁场\textbf{并不为零}，磁屏蔽效果远不如静电屏蔽；
  \item 改善办法：加厚罩壁或多层铁壳，把漏入的残余磁通逐次屏蔽；
  \item 不宜用于高频交变磁场，否则罩内产生很大的铁损（涡流、磁滞损耗）。
\end{itemize}

\corepoint{磁路计算三个易错点：(1) 把磁阻写成 $\mu S/l$（应为 $l/(\mu S)$，$\mu$ 对应 $\gamma$ 而非 $1/\gamma$）；(2) 漏算气隙——气隙磁阻常是铁心磁阻的几十上百倍，漏算会使 $NI$ 严重偏小；(3) 对铁磁段直接用 $H=B/\mu$ 而不查 $B$—$H$ 曲线（铁磁质 $\mu$ 不是常数）。}
